\clearpage
\subsubsection{多视角点云与后处理作用}
\ReportFigure[223bp]{raw_vs_fused.png}{49 个参考视角的点云对比。左图为所有预测深度直接变换到世界坐标系后的叠加，右图为阈值 0.8 的融合结果。两图使用相同观察方向与空间范围，未单独裁掉左图的外围点。}

直接将各视角点云放在一起后，建筑周围出现了较多杂乱点，部分点遮挡了主体。经过后处理，建筑外形更容易辨认，但屋顶和轮廓附近仍有少量残留。这说明，多视角信息不仅用于网络预测，也可以用来检查各张深度图之间是否互相支持。

置信度过滤先去掉模型判断不确定的像素；几何检查进一步去掉与其他视角不一致的点；深度平均将通过检查的观测结合起来；体素合并则减少同一区域的重复点。\textbf{点数的减少既包括不可靠点的去除，也包括重复点的合并，不能全部解释成错误点被删除。}在各视角共同判断错误时，后处理也可能保留错误结果。

\section{总结}
\subsection{对多视角重建的直观认识}

本次实验让我们看到，多张普通图像配合相机参数，已经能够恢复出较清楚的建筑点云。与实验1主要依靠单张图像中学到的场景经验不同，实验2还利用了不同观察位置之间的对应关系。已知相机内外参后，来自多个视角的信息能够为深度提供额外约束；重建尺度也随相机平移的尺度确定，不需要再用真实深度拟合一组尺度与偏移。

不过，输入视图数增加后，变化并不是每个位置都同样明显。建筑主体通常比较稳定，遮挡、背景和细小轮廓仍较难处理。我们也体会到，除了观察深度图是否平滑，还需要把结果变成点云，从空间中查看是否出现多余表面或缺口。

\subsection{参数选择与后续思考}

置信度阈值反映了保留范围与过滤强度之间的取舍：较低的阈值能保留更多点，但周围的杂点也更多；较高的阈值使结果更干净，同时可能去掉本来有用的观测。结合本次图像与点云展示，我们保留 0.8 作为默认设置，同时保存另外两组结果供比较，不将它视作所有场景都适用的最优值。

本次数据没有真实深度或参考点云，因此我们能判断的是可见形状与多视角一致性的变化，还不能给出绝对精度排名。后续如果能加入参考扫描，可以进一步比较重建误差和缺失程度，也可以改变拍摄位置，观察增加有用视角是否比单纯增加图像数量更有效。
